\vspace{-0.85cm}
\section{Introduction}
%\textcolor[RGB]{192,0,0}{\XSolidBrush}
% Table generated by Excel2LaTeX from sheet 'Sheet1'
\begin{table*}[!ht]
  \centering
  \caption{\textbf{Comparisons with Previous Works}. SceneScribe-1M is a large-scale video dataset with comprehensive geometric and semantic annotations. In the Geometric Annotation column, Depth map, Camera Pose, and 3D Tracks are abbreviated as D., C., and P., respectively.}
  \vspace{-0.2cm}
  \resizebox{\linewidth}{!}{
    \begin{tabular}{clcccccc}
    \toprule
    \multicolumn{1}{c}{\textbf{Type}} & \textbf{Dataset} & \textbf{Domain} & \textbf{Dynamic} & \textbf{Sem. Ann.} & \textbf{3D. Ann.} & \textbf{\#Scene Clips } & \textbf{\#Frames} \\
    \midrule
    \multirow{7}[2]{*}{3D Perception} & RealEstate10K~\cite{zhou2018stereo} & Indoor-Real & \XSolidBrush    & N/A & C.    & 80K   & 10M \\
    &BlendedMVS~\cite{yao2020blendedmvs} & Open-Synthetic & \XSolidBrush    & Single Label & D. C.  & 113   & 17K \\
    &CO3Dv2~\cite{reizenstein2021common} & Object-Real & \XSolidBrush    & Single Label & C.    & 19K   & 1.5M \\
    &PointOdyssey~\cite{zheng2023pointodyssey} & Object-Synthetic &  \CheckmarkBold    & N/A & D. C.  P. & 159   & 200K \\
    &CamVid-30K~\cite{zhao2024genxd} & Open-Real & \CheckmarkBold   & N/A & C.    & 30K   & - \\
    &Multi-Cam Video & Open-Synthetic & \CheckmarkBold   & Single Label & C.    & 136K  & 11M \\
    &DynPose-100K~\cite{rockwell2025dynamic} & Open-Real & \CheckmarkBold   & {Short Caption} & C.    & 100K  & 6.8M \\
    \midrule
    \multirow{3}[2]{*}{\makecell{Generation\\\& \\Understanding}} &HD-VILA-100M~\cite{xue2022advancing} & Open-Real & {\CheckmarkBold}   & Short Caption & N/A   & 103M  & 760k \\
    &Panda-70M~\cite{chen2024panda} & Open-Real & {\CheckmarkBold}   & Short Caption & N/A   & 70M   & 167K \\
    &Koala-36M~\cite{wang2025koala} & Open-Real & {\CheckmarkBold}   & Long Caption & N/A   & 36M   & 172k \\
    \midrule
    \multirow{3}[2]{*}{WFM}&Sekai-Real~\cite{li2025sekai} & Open-Real & {\CheckmarkBold}   & Structured Caption & D. C.  & $\sim$0.4M  & $\sim$40M \\
&SpatialVID~\cite{wang2025spatialvid} & Open-Real & {\CheckmarkBold}   & Structured Caption & D. C.  & $\sim$2M    & 123.6M \\
    &SceneScribe-1M (\textbf{ours}) & Open-Real & {\CheckmarkBold}   & Structured Caption & D. C.  P. & $\sim$1M    & 156.7M \\
    \bottomrule
    \end{tabular}}%
  \label{tab:source_video}%
  \vspace{-0.3cm}
\end{table*}%



\label{sec:intro}
%spatial intelligence
In recent years, the rapid advancement of 3D geometric perception and video synthesis have significantly accelerated research in world foundation models (WFMs)~\cite{deepmind_genie3_2024,ali2025world,li2025sekai}. 
Collectively, these technologies enable WFMs to perceive, simulate, and interact effectively within dynamic environments.
Such capabilities integrated by WFMs are critical for promoting transformative developments in areas such as augmented reality~\cite{huang2025voyager}, robotics~\cite{fu2025learning,cen2025worldvla}, and autonomous driving~\cite{li2025omninwm,li2025drivevla}.
However, the scarcity of sufficiently large and high-quality datasets restricts the potential of existing models in both 3D perception and video synthesis, thereby further hindering the prospects of WFMs.
% Such limitations pose significant challenges for practical applications demanding precise geometric details and high visual fidelity.

 






Current efforts to address data challenges related to 3D perception can be categorized into two main paradigms. 
%1synthetic data,
One common strategy~\cite{cabon2020virtual,yao2020blendedmvs,bai2025recammaster} follows a data synthesis pipeline within virtual engines, automatically generating ground-truth camera poses and corresponding geometric annotations. 
Nevertheless, these approaches introduce a domain gap and overlook complex physical interactions.
%2semi-automatic annotation
Alternatively, another prevalent routine attempts to efficiently annotate real-world data by SfM~\cite{schonberger2016structure} or SLAM~\cite{mur2015orb} systems. Apart from the sparsity of camera trajectory annotations in static scenes~\cite{rockwell2025dynamic}, the annotation scale and diversity for dynamic scenes are also limited by computational overhead~\cite{zhou2018stereo,zhao2024genxd}.
Beyond 3D perception, video generation data with rich semantic information is also essential for building WFMs. Notably, current open-world datasets~\cite{nan2024openvid,wang2025koala,chen2024panda} have somewhat alleviated the issues of limited data and annotation scarcity present in previous studies~\cite{soomro2012ucf101,zhu2022celebv,yuan2024chronomagic}.
Nonetheless, since these datasets are tailored for video generation (e.g., text-to-video~\cite{kong2024hunyuanvideo}), they lack geometric annotations, consequently leaving the semantic and motion diversity required by WFMs insufficiently examined.
Despite the above progress of single-modal datasets, advances in WFMs remain fundamentally constrained by the inadequacy of large-scale datasets that comprehensively capture 3D geometric and fine-grained semantic properties.






In this paper, we introduce SceneScribe-1M, a large-scale, multi-modal video dataset that facilitates the critical intersection of 3D geometric perception and video synthesis (as shown in Figure~\ref{teaser}).
By incorporating powerful models in proprietary domains (\textit{i.e.}, Qwen2.5-VL-72B~\cite{bai2025qwen2}, MegaSaM~\cite{li2025megasam}, and TAPIP3D~\cite{zhang2025tapip3d}), we deploy over 1,000 GPUs to implement our annotation pipeline on large-scale videos.
SceneScribe-1M comprises one million in-the-wild videos, amounting to over 4,000 hours, each extensively annotated with detailed textual descriptions, precise camera parameters, continuous video depths, and consistent 3D point tracks.
Crucially, our curation establishes criteria across four key aspects, informed by both semantic and geometric annotations: video parameters, semantic information, camera motion, and object motion. Raw videos are meticulously examined based on these indicators to ensure content diversity and motion richness.
% This subset is derived through a multi-view filtering mechanism that flexibly controls the inclusion of dynamic objects.
We further devise a filtering mechanism for SceneScribe-MVS subset construction, aiming to accommodate multi-view tasks that prefer static objects. 
This filter disentangles the camera and object motion, controlling the dynamic object inclusion without compromising camera motion intensity.
To establish rigorous benchmarks, we leverage SceneScribe-1M for core 3D perception, including monocular depth estimation, scene reconstruction, and dynamic point tracking. 
Moreover, SceneScribe-1M serves as a pivotal resource for advancing generative tasks such as text/pose-to-video synthesis, supporting precise view control over camera motion.

%which are vital for creating spatially coherent and controllable visual sequences.

\begin{figure*}[!ht]
  \centering
    \includegraphics[width=\linewidth]{figures/framework.pdf}
      \vspace{-0.6cm}
  \caption{\textbf{Curation Pipeline} for SceneScribe-1M consist of: (a) We begin by collecting large-scale videos from various sources; (b) Raw videos undergo specification and content inspection, with temporal segmentation models employed to ensure continuity; and (c) We integrate Qwen2.5-VL-72B~\cite{bai2025qwen2}, MegaSaM~\cite{li2025megasam}, and TAPIP3D~\cite{zhang2025tapip3d} to perform comprehensive geometric and semantic annotations.}
  \vspace{-0.2cm}
  \label{fig:framework}
\end{figure*}

In summary, our primary contributions are as follows:
\begin{itemize}
    \item \textbf{Comprehensive Video Annotations}: SceneScribe-1M contains over 4,000 hours of video data, accompanied by essential geometric and semantic annotations. These annotations provide a unified resource that facilitates both large-scale 3D perception and video generative tasks.
    \item \textbf{Curated Videos with Semantic and Motion Diversity}: 
     SceneScribe-1M is curated with semantic and geometric indicators for content and motion diversity. 
    We also introduce a multi-view filter for SceneScribe-MVS to limit dynamic objects while preserving camera motion.
    \item \textbf{Extensive Downstream Evaluation}: The potential versatility of SceneScribe-1M is demonstrated by its applicability across diverse downstream tasks, including 3D geometric perception and video synthesis, which in turn highlight both the effectiveness and the quality of the dataset.
\end{itemize}





% %-------------------------------------------------------------------------
% \subsection{Language}



